\section{文献综述与理论基础}

\subsection{供应链风险测度研究}

供应链风险研究通常强调风险并非只来自单一断供事件，而是来自需求波动、供应中断、政策限制、物流受阻和关键节点失效等多类冲击的叠加 \citep{chopra2004managing,juttner2005supply}。对进口依赖型产业而言，风险暴露不仅表现为某一来源国份额较高，也体现为供应商结构过度集中、替代来源不足以及来源结构在短期内剧烈波动。因而，供应链安全评估需要同时观察脆弱性和恢复能力：前者对应冲击发生时受影响的可能程度，后者对应通过多元来源和替代路径恢复供给的能力。

这一视角为本文 SIRI 指数提供理论基础。本文不把半导体进口风险简化为“美国份额”单一指标，而是将来源集中、政策暴露、替代不足和结构波动同时纳入评价。这样的设定有助于识别两类容易被忽略的情形：一是对美份额不高但非美来源高度集中的结构性风险，二是总进口规模稳定但来源替代能力不足导致的隐性脆弱性。

\subsection{半导体贸易与出口管制研究}

半导体产业具有高度全球分工特征，设计、制造、设备、材料、封装测试和终端需求分布在不同经济体之间。全球价值链研究指出，信息技术降低了跨境协调成本，使生产环节更容易在国家间分解和重组 \citep{baldwin2016great}。因此，半导体贸易结构不能只看总体进口额，更需要在产品级层面区分制造设备、存储器、处理器及其他集成电路等类别，因为不同产品对应的来源国组合、技术约束和替代空间并不相同。

近年来，美国针对先进计算芯片、半导体制造设备和相关技术的出口管制构成了本文的重要政策背景 \citep{us_bis_controls}。但从贸易数据解释角度看，进口来源变化可能同时受到出口管制、全球半导体周期、中国需求结构、第三方产能扩张和企业库存调整等因素影响。因此，本文将出口管制作为观察结构演化的关键时间背景，而不直接把产品级贸易份额变化表述为单一政策冲击的因果结果。

\subsection{贸易网络与风险预测研究}

网络视角能够把贸易关系从双边份额扩展为节点、边和权重共同构成的结构对象。社会与经济网络理论表明，节点中心性、连接密度和替代路径会影响冲击传播和系统韧性 \citep{jackson2010social}。用于中国半导体相关产品进口研究时，年度贸易网络可以同时刻画来源集中度、主要来源国更替、非美替代路径以及某一关键来源受限时的潜在传播影响。

在预测层面，图表示学习为利用网络结构生成预测输入提供了技术工具，例如 GraphSAGE 将节点邻域信息聚合为可用于下游任务的表示 \citep{hamilton2017inductive}。但这类方法在本文中只构成预测扩展的输入视角，不能替代因果识别。也就是说，贸易网络特征可以用于探索下一年 SIRI 风险是否具有可预测性，但不能据此断言某一政策事件造成了特定贸易结构变化。

\subsection{文献评述与本文定位}

综合来看，既有研究为供应链风险来源、全球价值链分工和网络结构分析提供了基础，但直接服务于本文问题时仍有三点需要补足：第一，半导体相关产品内部差异明显，使用总体贸易额容易掩盖设备类和集成电路类产品的风险来源差异；第二，单一来源国份额难以同时反映集中、替代不足和结构波动；第三，预测模型若缺少明确证据边界，容易把相关性或预测性能误写成政策因果效应。

据此，本文定位为产品级 BACI 数据驱动的进口来源结构研究：先基于 HS6 产品--年份面板刻画中国半导体相关产品进口来源演化，再构建可解释的 SIRI 风险指数衡量来源集中、政策暴露、替代不足和结构波动，随后用固定效应回归给出谨慎的统计证据边界，最后开展 BACI-only 贸易网络下一年 SIRI 预测扩展。该定位强调可复现的数据链条和可解释的风险评价，同时明确预测扩展不替代前文的描述统计与回归检验。
